% This Source Code Form is subject to the terms of the Mozilla Public
% License, v. 2.0. If a copy of the MPL was not distributed with this
% file, You can obtain one at https://mozilla.org/MPL/2.0/.

\documentclass[a4paper, 12pt]{article}

\usepackage[english, russian]{babel}
\usepackage[T2A]{fontenc}
\usepackage[utf8]{inputenc}
\setlength{\parindent}{1cm}
\usepackage{graphicx}
\usepackage{amsmath}
\usepackage{amssymb}
\usepackage{amsfonts}
\usepackage{array}
\usepackage{booktabs}
\usepackage{wrapfig}
\usepackage{caption} %заголовки плавающих объектов
\captionsetup[table]{justification=centering} %установки для заголовков таблиц
\captionsetup[figure]{justification=centering} %установки для заголовков таблиц

\begin{document}
	\begin{center}
		\subsubsection*{Расчётное задание по математической физике}
		\subsubsection*{Глава №2: Задачи Штурма-Лиувилля  в  математической  физике}
		\subsubsection*{Выполнил студент третьего курса Физико-механического института, высшей школы ФФИ: Куликовских Илья}
		\subsubsection*{Группа: 5030302/30801}
		\subsubsection*{Вариант 6}
		\subsubsection*{Преподаватель: Елатонцев Вадим Александрович}
	\end{center}
	
	\quad \textbf{Задание 2.01.07} Показать что уравнение:
	
	\begin{equation}
		\dfrac{d^2 y}{dx^2} + \dfrac{2}{x} \dfrac{dy}{dx} + y = 0
	\end{equation}
	
	имеет решение $y_1(x) = \dfrac{1}{x} \sin x$ . Найти второе линейно-независимое решение $y_2 (x)$ по формуле:
	
	\begin{equation}
		y_2 (x) = y_1 (x) \displaystyle\int\dfrac{e^{-A_1 (x)}}{y_1^2 (x)} dx 
	\end{equation}
	
	\quad где $A_1 = \displaystyle\int \dfrac{2}{x} dx$ .  Вычислить определитель вронского $W(x)$.
	
	\quad \textbf{Решение.} Проверим, что $y = \dfrac{1}{x}\sin x$ действителььно является решением ОЛДУ (1):
	
	\begin{center}
		$\dfrac{2}{x^3}\sin x -\dfrac{2}{x^2}\cos x - \dfrac{1}{x}\sin x - \dfrac{2}{x^3}\sin x + \dfrac{2}{x^2}\cos x + \dfrac{1}{x}\sin x = 0 $ 
	\end{center}
	
	\quad Далее определим второе решение ОЛДУ (1) по формуле (2) и проверим его линейную зависимость с первым, вычислив определитель Вронского $W(x)$:
	
	\begin{center}
		$\displaystyle\int\dfrac{2}{x}dx = 2\ln|x| \implies e^{-2\ln|x|} = \dfrac{1}{x^2}$

		$y_2(x) = \dfrac{1}{x}\sin x \displaystyle\int\dfrac{1}{x^2 \frac{1}{x^2} \sin^2 x} dx = \dfrac{1}{x}\sin x \ctg x = \dfrac{1}{x}\cos x $
	\end{center}
	
	\quad Проверим получившиеся решение: 
	
	\begin{center} 
		$\dfrac{2}{x^3} \cos x + \dfrac{2}{x^2}\sin x - \dfrac{1}{x}\cos x - \dfrac{2}{x^3}\cos x - \dfrac{2}{x^2}\sin x + \dfrac{1}{x}\cos x = 0 $
	\end{center}
	
	\quad Вычислим определитель Вронского:
	
	\begin{equation}
		W(x) = \left|\begin{array}{cc}
			\dfrac{1}{x}\sin x & \dfrac{1}{x}\cos x \\
			\dfrac{1}{x}\cos x - \dfrac{1}{x^2} \sin x & -\dfrac{1}{x}\sin x - \dfrac{1}{x^2}\cos x \\
		\end{array}\right|
	\end{equation}
	
	\begin{equation}
		W(x) = -\dfrac{1}{x^2}\sin^2 x - \dfrac{1}{2x^3}\sin2x - \dfrac{1}{x^2}\cos^2 x + \dfrac{1}{2x^3}\sin 2x = -\dfrac{1}{x^2}\implies 
	\end{equation}
	
	\quad $\implies $ Решения являются линейно независимыми и следовательно образуют фундаментальную систему решений (ФСР) ОЛДУ (1) $\blacksquare$.
		
	\quad \textbf{Задание 2.01.08} Найти общее решение неоднородного линейного дифференциального уравнения методом вариации постоянных: 
	
	\begin{equation}
		\dfrac{d^2 y}{dx^2} + \dfrac{2}{x} \dfrac{dy}{dx} + y = \dfrac{ \ctg x }{x}
	\end{equation} 
	
	\quad \textbf{Решение.} Общее решение НЛДУ является суммой общего решения ОЛДУ и частоного решения этого НЛДУ. Для выражения общего решения этого ОЛДУ воспользуемся результатои прошлой задачи:
	
	\begin{center}
		$\dfrac{d^2 y}{dx^2} + \dfrac{2}{x} \dfrac{dy}{dx} + y = 0$
	\end{center}
	\begin{equation}
		y(x) = C_1 y_1 (x) + C_2 y_2 (x) = \dfrac{C_1}{x}\sin x + \dfrac{C_2}{x}\cos x
	\end{equation}
	
	\quad Частное решение уравнения (5) имеет вид: 
	
	\begin{equation}
		y_{\text{част}}(x) = C_1 (x) y_1 (x) + C_2 (x) y_2 (x) 
	\end{equation}

	\quad Для определения констант $C_1 (x)$, $C_2 (x)$ воспользуемся методом вариации постоянных. Рассмотрим систему: 
	
	\begin{equation}
		\begin{cases}
			y_1 (x)\dfrac{dC_1}{dx} + y_2 (x)\dfrac{dC_2}{dx} = 0 \\
			\\
			\dfrac{dy_1}{dx} \dfrac{dC_1}{dx} + \dfrac{dy_2}{dx}\dfrac{dC_2}{dx} = \dfrac{\ctg x}{x} 
		\end{cases} \implies \left[\begin{array}{cc|c }
			y_1 & y_2 & 0 \\
			& & \\
			\dfrac{dy_1}{dx} & \dfrac{d y_2}{dx} & \dfrac{\ctg x}{x}
		\end{array}\right]
	\end{equation}
	
	\begin{center}
		$\implies \Delta = W(x) = -\dfrac{1}{x^2}$; $\Delta_1 = -\dfrac{\ctg x}{x}y_2(x)$; $\Delta_2 = \dfrac{\ctg x}{x} y_1 (x)$
	\end{center}
	
	\quad $\dfrac{dC_1}{dx} = \dfrac{\Delta_1}{\Delta} = \dfrac{x^2 \ctg x}{x^2}\cos x = \dfrac{\cos^2 x}{\sin x} = \dfrac{1}{\sin x} - \sin x \implies $
	
	$\implies C_1 = \displaystyle\int\left(\dfrac{1}{\sin x} -\sin x\right)dx = \displaystyle\int\dfrac{dx}{\sin x} - \displaystyle\int\sin x dx = \displaystyle\int\dfrac{2(1 + t^2)}{2t(1 + t^2)}dt + +\cos x = \ln|\tg \dfrac{x}{2}| + \cos x + \tilde{C_1}$
	
	\quad $\dfrac{dC_2}{dx} = \dfrac{\Delta_2}{\Delta} = -\dfrac{x^2 \cos x}{x^2} = -\cos x \implies  C_2 = -\displaystyle\int\cos x dx =\newline= -\sin x + \tilde{C_2}$
	
	\quad Таким образом частное решение НЛДУ (5) имеет вид:
	
	\begin{equation}
		y_{\text{част}}(x) = \dfrac{1}{x}\left( \ln|\tg\dfrac{x}{2}| + \cos x \right)\sin x - \dfrac{1}{x}\sin x \cos x
	\end{equation}
	
	\quad Общее решение НЛДУ (5) таким образом равно:
	
	\begin{equation}
		y_{\text{общ}} = \dfrac{C_1}{x}\sin x + \dfrac{C_2}{x}\cos x + \dfrac{1}{x}\left( \ln|\tg\dfrac{x}{2}| + \cos x \right)\sin x - \dfrac{1}{x}\sin x \cos x
	\end{equation} 
	
	\quad \textbf{Ответ: } $y_{\text{общ}} = \dfrac{C_1}{x}\sin x + \dfrac{C_2}{x}\cos x + \dfrac{1}{x}\left( \ln|\tg\dfrac{x}{2}| + \cos x \right)\sin x - -\dfrac{1}{x}\sin x \cos x$
	
	\quad \textbf{Задание 2.02.04} Показать что уравнение: 
	
	\begin{equation}
		y\dfrac{\partial U}{\partial x} - x \dfrac{\partial U}{\partial y} = 0
	\end{equation}
	
	имеет общее решение $U(x, y) = \varphi(\sqrt{x^2 + y^2}) $. Использовать метод характеристик. 
	
	\quad \textbf{Решение.} Распишем уравнения метода характеристик
	
	\begin{center}
		$\dfrac{dy}{dt} = -x$; $\dfrac{dx}{dt} = y \implies $
		
		$\implies \dfrac{dy}{x} = -\dfrac{dx}{y} \implies y^2 + x^2 = (\sqrt{x^2 + y^2})^2 = Const$
		
	\end{center}
	
	\quad Таким образом, $U(x, y) = f(Const) = f(\sqrt{x^2 + y^2}) $ $\blacksquare$.
	
\end{document}